\clearpage
\VerdierPhysicalPageAnchor{0218}%
\begingroup\small\setlength{\parskip}{0pt}\setlength{\emergencystretch}{1.5em}
\setlength{\abovedisplayskip}{4pt plus 2pt minus 2pt}%
\setlength{\belowdisplayskip}{4pt plus 2pt minus 2pt}%

\VerdierSemanticAnchor{ch3d-3-3-1-2}
\noindent\textit{exactes $S_i^1*f_{i-1}$, $1<i\leq n$,
$f_i*{S''}_i^1$, $1\leq i<n$, ${S'}_i^1*g_{i-1}$, $1<i\leq n$,
$g_i*{S''}_i^1$, $1\leq i<n$, soient définies ((3.2.8.1) et (3.2.8.2))
et tels que :}

\begin{equation}
\left\{
\begin{aligned}
\delta^1({S''}_n^1)
  &=\delta^1(S_n^1*f_{n-1}),\\
\delta^1({S''}_n^1)
  &=\delta^1({S'}_n^1*g_{n-1}),\\
\delta^1(f_i*{S''}_i^1)
  &=\delta^1(S_i^1*f_{i-1}),&&1<i\leq n-1,\\
\delta^1(g_i*{S''}_i^1)
  &=\delta^1({S'}_i^1*g_{i-1}),&&1<i\leq n-1,\\
\delta^1(f_1*{S''}_1^1)
  &=\delta^1(S_1^1),\\
\delta^1(g_1*{S''}_1^1)
  &=\delta^1({S'}_1^1).
\end{aligned}
\right.
\tag{3.3.1.2}
\end{equation}

\VerdierSemanticAnchor{ch3d-3-3-1-proof}
En effet, l'assertion a) résulte immédiatement de (3.2.2) et (3.2.5) et
l'assertion b) résulte de la description de la relation d'équivalence induite
par l'application
$E_{\mathcal A}^{n}(X,Y)\longrightarrow
\operatorname{Ext}_{\mathcal A}^{n}(X,Y)$ ((3.2.2) et (3.2.9)).

\VerdierSemanticAnchor{ch3d-3-3-2}
\medskip\noindent\textbf{3.3.2.}\quad
Rappelons (1.2.11) que la catégorie $\mathcal A^{\operatorname{Ext}}$ est la
sous-catégorie pleine de $D(\mathcal A)$ définie par les translatés des
objets de $\mathcal A$. Soit
$F\colon\mathcal A^{\operatorname{Ext}}\longrightarrow\mathcal C$ un
foncteur (à valeurs dans une catégorie non nécessairement additive).
Désignons, pour tout entier $n$, par
$F^n\colon\mathcal A\longrightarrow\mathcal C$ le foncteur
$X\longmapsto F(X[n])$. De même, pour toute suite exacte de $\mathcal A$,
$S^1=0\longrightarrow X\longrightarrow Y\longrightarrow Z
\longrightarrow0$, désignons par
$\delta_F^n(S^1)\colon F^n(Z)\longrightarrow F^{n+1}(X)$ le morphisme

\VerdierSemanticAnchor{ch3d-3-3-2-delta-definition}
\[
F\bigl(\delta^1(S^1)[n]\bigr)
\colon F(Z[n])\longrightarrow F(X[n+1]).
\]

\VerdierSemanticAnchor{ch3d-3-3-2-sfl}
La suite des foncteurs $F^n$, $n\in\mathbb Z$, munie des morphismes
$\delta_F^\ast$ possède les propriétés suivantes (3.2.9) :

\noindent\textbf{SFL I)} Tout diagramme commutatif :

\VerdierSemanticAnchor{ch3d-3-3-2-sfl-i-source-diagram}
\begin{center}
\resizebox{0.86\textwidth}{!}{%
$\begin{tikzcd}[ampersand replacement=\&,column sep=3.0em,row sep=2.8em]
S=0 \arrow[r]
  \& X \arrow[r] \arrow[d,"f"']
  \& Y \arrow[r] \arrow[d]
  \& Z \arrow[r] \arrow[d,"g"]
  \& 0 \\
{S'}=0 \arrow[r]
  \& X' \arrow[r]
  \& Y' \arrow[r]
  \& Z' \arrow[r]
  \& 0
\end{tikzcd}$%
}
\end{center}
\endgroup
